% !PRINTED 669
\weberpage{669}
\begin{center}
§ 190. \gesperrt{Verzweigungspunkte und Verzweigungszahlen.}
\end{center}

Nach § 184 gibt es nur eine endliche Anzahl von Punkten, in denen $z-c$ Null in höherer als der ersten Ordnung wird. Diese heißen die Verzweigungspunkte von $\Omega$ in $z$.

Wir konstruieren ein Polygon $\mathfrak{Z}_z$, das $(e-1)$mal jeden Punkt $\mathfrak{P}$ enthält, in dem $z-c$ oder $1:z$ Null in der $e$ten Ordnung wird. Dieses Polygon heißt das \gesperrt{Verzweigungspolygon von $\Omega$ in bezug auf $z$}, und seine Ordnung

\weq{2}{w_z = \Sigma(e-1)}

heißt die \gesperrt{Verzweigungszahl in bezug auf $z$}\footnote{Breitet man die \gesperrt{Riemann}sche Fläche über der $z$-Ebene aus, so ist $n$ die Anzahl der übereinanderliegenden Blätter, $w_z$ ist die Anzahl der Verzweigungspunkte dieser Fläche.}.

In § 184, (20) haben wir die Formel erhalten:

\weq{3}{N(t-\tau) \equiv (t-\tau_1)(t-\tau_2)\ldots(t-\tau_n) \ (\mathrm{mod}\ z-c),}

worin $\tau$ die Basisform des Körpers $\Omega$ nach $z$ bedeutet und $\tau_1, \tau_2, \ldots, \tau_n$ die gleichen oder verschiedenen konstanten Funktionale, denen $\tau$ nach den Primfaktoren von $z-c$ kongruent wird, also die Werte, die $\tau$ in den nach $z$ konjugierten Punkten $\mathfrak{P}_1, \mathfrak{P}_2, \ldots, \mathfrak{P}_n$ annimmt. In jedem dieser Punkte muß (3) in eine richtige Gleichung übergehen, und daraus folgen z.\,B.:

\weq{4}{N(\tau) = \tau_1 \tau_2 \ldots \tau_n \,,}

\weq{5}{S(\tau) = \tau_1 + \tau_2 + \cdots + \tau_n}

als Werte der rationalen Funktionale $N(\tau)$ und $S(\tau)$ für $z=c$. Setzt man für die Funktionalvariablen rationale Funktionen von $z$, so gelten entsprechende Formeln für die Funktionen in $\Omega$, vorausgesetzt, daß nicht einer der Ausnahmefälle $0\,.\,\infty$, $\infty+\infty$ eintritt.

Ist daher $\eta_1, \eta_2, \ldots, \eta_n$ eine Basis von $\Omega$, und sind $\eta_{i,1}, \eta_{i,2}, \ldots, \eta_{i,n}$ konjugierte Werte nach $z$ von $\eta_i$, so ist nach (5) und der Definition der Diskriminante in § 173

\weq{6}{\Delta(\eta_1, \eta_2, \ldots, \eta_n) = (\Sigma \pm \eta_{1,1}\eta_{2,2}\ldots\eta_{n,n})^2.}

Stehen zwei Funktionen $z$ und $z'$ in $\Omega$ in einer linearen Beziehung:

\weq{7}{z' = \dfrac{az+b}{cz+d},}
%CONT>>
